% C.4ab. Fully audited criterion and method exclusion; no priority claim.
\paragraph{C.4ab (Proved: a partial Fourier test for the surface case).}

\emph{Status: proved criterion and method exclusion. This is an
equivalent formulation of the residual splitting
problem, not a solution of that problem or a claim of priority.}

Let $A$ be a complex abelian surface with a primitive polarization
$\ell=c_1(L)$ of type $(1,2)$, and assume
\[
 \operatorname{NS}(A)=\mathbb Z\ell,\qquad
 \operatorname{End}_{\mathrm{HS}}H^1(A,\mathbb Q)=\mathbb Q,
 \qquad
 \operatorname{End}_{\mathrm{HS}}T(A)_{\mathbb Q}=\mathbb Q,
\]
Here we put $T(A)_{\mathbb Q}=\ell^\perp\subset H^2(A,\mathbb Q)$.
The last condition is an explicit assumption. We do not infer it
from a claim about a special surface.

These hypotheses include the Mumford--Tate general case, in the
sense of \citet[p.~2473]{cyVoisin2024}. Both relevant symplectic
representations are absolutely irreducible: the standard one
and the primitive part of its second exterior power. Put
\[
 X=A\times\widehat A,\qquad P=c_1(\mathcal P_A),\qquad
 \ell_1=\operatorname{pr}_A^*\ell,
 \quad \ell_2=\operatorname{pr}_{\widehat A}^*\widehat\ell,
\]
where the primitive dual polarization satisfies
$\lambda_{\widehat\ell}\lambda_\ell=[2]_A$.

For an integral cycle $Z\in\mathrm{CH}^2(A\times A)$, we use
the convention
\[
 Z^*\alpha=(p_1)_*\bigl([Z]\smile p_2^*\alpha\bigr).
\]
There are integers $t,s,n$ such that $Z^*$ acts by $t$ on
$H^1(A,\mathbb Q)$, by $s$ on $T(A)_{\mathbb Q}$, and by $n$
on $\mathbb Q\ell$. The following conditions are equivalent:
\begin{enumerate}
\item $A$ is an actual direct summand of a product of Jacobians;
\item some integral $Z\in\mathrm{CH}^2(A\times A)$ has
      $t-s$ odd.
\end{enumerate}
More precisely, for every such $Z$ we construct an integral curve
cycle on $X$ whose mixed Matsusaka endomorphism is
\[
 m_Z\operatorname{id}_A,
 \qquad
 m_Z=3(t^2-s)-\frac{n-s}{2},
 \qquad m_Z\equiv t-s\pmod2.                         \tag{PF.1}
\]
The integer $m_Z$ is an actual integral coefficient, not the
coefficient of a rationally divided Chern character.

\emph{Step 1: the integral gluing condition.}
The action of an integral correspondence preserves integral
cohomology. The Hodge-endomorphism assumptions therefore make
$t,s,n$ integers: each scalar preserves a nonzero free lattice,
and $\ell$ is primitive. There are no rational Hodge homomorphisms
between $T(A)_{\mathbb Q}$ and $\mathbb Q\ell$. Indeed a map in
one direction would give a rational Hodge class in $T(A)$, and
duality gives the same conclusion in the other direction.

The lattice $H^2(A,\mathbb Z)$ is unimodular, and $\ell^2=4$.
Choose $w$ in this lattice with $(w,\ell)=1$. The decomposition
of $Z^*$ into its two rational eigenspaces gives
\[
       Z^*w-sw=\frac{n-s}{4}\ell.
\]
The left side is integral, and $\ell$ is primitive. Consequently
\[
                         n-s\in4\mathbb Z.            \tag{PF.2}
\]
Write $\Delta_{22}$ for the $(2,2)$ component of the diagonal.
The same component of $[Z]$ is
\[
                  s\Delta_{22}+\frac{n-s}{4}\ell\boxtimes\ell.
                                                               \tag{PF.3}
\]
Its $(1,3)$ component is $t\Delta_{13}$.

\emph{Step 2: an integral Fourier construction.}
Write $Z=\sum_i m_iZ_i$ and choose resolutions
$f_i:\widetilde Z_i\to Z_i\hookrightarrow A\times A$.
The integral class
\[
                  K_Z=\sum_i m_i(f_i)_*[\mathcal O_{\widetilde Z_i}]
                         \in K_0(A\times A)
\]
has $\operatorname{ch}_0(K_Z)=\operatorname{ch}_1(K_Z)=0$
and $\operatorname{ch}_2(K_Z)=[Z]$. This follows from the
leading-degree part of Grothendieck--Riemann--Roch; the cycle
pushforward has coefficient one because each resolution is
birational onto $Z_i$.

Apply the integral Fourier--Mukai functor in the second factor:
\[
 E_Z=R(q_{a,y})_*
       \bigl(p_{a,x}^*K_Z\otimes p_{x,y}^*\mathcal P_A\bigr)
       \in K_0(A\times\widehat A).
                                                               \tag{PF.4}
\]
Here the three factors before pushforward are $A_a,A_x,\widehat A_y$,
and the Poincar\'e bundle in (PF.4) is on $A_x\times\widehat A_y$.
We use the pushforward of perfect complexes, so no WIT condition
or assertion that $E_Z$ is a single vector bundle is needed.
Its Chern classes are integral Chow classes, also when $E_Z$
is virtual. The integral functor and its cohomological transform
are those of \citet[\S1, (1.4), and Theorem 2.2]{cyMukai1981PF};
the Riemann--Roch formula is explicitly written in
\citet[(2.6)--(2.8), author p.~5]{cyYoshioka2001PF}.
The tangent bundle of an abelian variety is trivial, so the
cohomological transform is
\[
                   \mathcal F_x(u)=(q_{a,y})_*
                                  \bigl(p_{a,x}^*u\,e^{P_{x,y}}\bigr).
\]

We fix the usual normalized Poincar\'e sign, for which
\[
 \mathcal F_x(\Delta_{13})=P,\qquad
 \mathcal F_x(\Delta_{22})=P^2/2,\qquad
 \mathcal F_x(\ell_x)=-\widehat\ell_y.
                                                               \tag{PF.5}
\]
These identities can also be checked directly in the exterior
algebra on $H^1(A_a)\oplus H^1(A_x)\oplus H^1(\widehat A_y)$.
For an integral basis, write
\[
 \ell_a=a_0a_1+2a_2a_3,\qquad
 \widehat\ell_y=2y_0y_1+y_2y_3,\qquad
 P_{a,y}=\sum_{j=0}^3a_jy_j.
\]
Multiplying $\prod_{j=0}^3(x_j-a_j)$ by
$\exp(\sum_jx_jy_j)$ and integrating the four $x_j$ gives
$\exp(P_{a,y})$, including every sign in (PF.5).

The Fourier transform takes bidegree $(a,b)$ to $(a,4-b)$.
Since $K_Z$ has no Chern-character components of degrees zero
or one, this gives
\[
                c_1(E_Z)=a\ell_1+tP\quad (a\in\mathbb Z).
\]
In particular there is no $\ell_2$ term: it could only come
from the vanished $(0,2)$ component of $\operatorname{ch}_1(K_Z)$.
The $(2,2)$ component of $\operatorname{ch}_2(E_Z)$ is, by
(PF.3)--(PF.5),
\[
                s\frac{P^2}{2}-\frac{n-s}{4}\ell_1\ell_2.
\]
Therefore the $(2,2)$ component of $c_2(E_Z)$ is
\[
               \frac{t^2-s}{2}P^2+\frac{n-s}{4}\ell_1\ell_2.
                                                               \tag{PF.6}
\]
All the omitted components have different bidegrees; after
intersection with $P$, they cannot contribute to
$H^3(A)\otimes H^3(\widehat A)$, which is the component controlling
the mixed endomorphism.

\emph{Step 3: the mixed integer.}
In the same exterior basis we have
\[
                         P^3=-3\ell_1\ell_2P.          \tag{PF.7}
\]
The integral cohomology class $P^3/6$ has mixed Matsusaka
endomorphism $\operatorname{id}_A$. This is the unimodular
Poincar\'e normalization already used in C.4i and C.4r; it
does not assert that $P^3/6$ is an integral Chow class.
For the actual integral curve
\[
                              z_Z=c_2(E_Z)\cdot P
                                       \in\mathrm{CH}^3(X),
\]
(PF.6)--(PF.7) give (PF.1). By (PF.2), its second summand is
even, while $t^2\equiv t\pmod2$.
Each normalized curve component of $z_Z$ supplies a factorization
of its mixed endomorphism through its Jacobian. Thus
$m_Z\in I_J(A)$. If $t-s$ is odd, $m_Z$ is odd. C.4n gives
$2\in I_J(A)$, and B\'ezout gives $1\in I_J(A)$, proving (1).

\emph{Step 4: the converse is integral.}
Suppose that there are curves $C_i$ and homomorphisms
\[
       u_i:A\to J(C_i),\qquad v_i:J(C_i)\to A,
                    \qquad \sum_i v_i u_i=\operatorname{id}_A.
\]
Choose a point on each curve and the corresponding Abel map
$a_i:C_i\to J(C_i)$. Pull the normalized Poincar\'e line bundle
on $J(C_i)\times C_i$ back by $u_i\times\operatorname{id}_{C_i}$;
let its integral divisor class be $D_i\in\mathrm{CH}^1(A\times C_i)$.
Set
\[
               Z=\sum_i
                  (\operatorname{id}_A\times(v_i a_i))_*D_i
                         \in\mathrm{CH}^2(A\times A).
\]
The action on $H^1$ is
$\sum_i u_i^*v_i^*=(\sum_i v_i u_i)^*=\operatorname{id}$.
The action on $T(A)$ factors through $H^2(C_i,\mathbb Q)$,
a Tate Hodge structure, and is zero. Hence $t=1,s=0$.
No division, isogeny inverse, or saturation of a cycle group
enters this construction. This proves the equivalence.

In this precise surface scope the conditions are also equivalent
to integral algebraicity of the first K\"unneth projector
$\delta_1$, namely ``splitness'' in
\citet[Definition 1.9, p.~2473]{cyVoisin2024}.
Indeed, such a projector has $t=1,s=0$, so the forward
construction applies. Conversely the normalized divisor $D_i$
above has only the component
$H^1(A)\otimes H^1(C_i)$. Its pushforward therefore has only
component $(1,3)$, and the sum acts as the identity on $H^1$.
Thus its cohomology class is exactly $\delta_1$.
We assert this equivalence only under the stated surface
hypotheses, not for abelian varieties of arbitrary dimension.

\emph{Step 5: the known residual fibre-product surface fails this test.}
Suppose $\pi:S\to A$ is a composition of blowups and
$\gamma:S\to Y$ is a generically finite morphism to a smooth
projective surface whose $H^1$ is zero and whose $H^2$ is
generated by divisor classes. Assume that $\gamma$ is finite
and generically separable. The cycle of its fibre product is
\[
             [S\times_Y S]=[\Delta_S]+[W],
\]
where $W$ is the residual correspondence, with its natural
cycle multiplicities. Its action is
$\gamma^*\gamma_*-\operatorname{id}$.
The first term is zero on $H^1(S)$ and on $T(S)$.
Pushing $W$ to $A\times A$ therefore gives $t=s=-1$,
so the mixed integer of its partial Fourier construction is
even. The same conclusion holds for integral sums of these
correspondences, multiplication graphs, and exterior products
of divisors: the parity $t-s$ is additive and vanishes on each.

This applies directly to the Galois-closure surface of
\citet[Proposition 4.5 and \S5, (5.1), author pp.~16--18]
{frC4BastianelliPirolaStoppino2010}: their resolved abelian
surface maps finitely with degree three to the rational ruled
surface $\mathbb F_3$, and their surface is the residual
fibre product. Its full-span Albanese property alone would not
justify this Fourier exclusion; the residual-correspondence
calculation does. We make no assertion here about every
codimension-two cycle on $A\times A$ or every Fourier transform
of an arbitrary sheaf.

\emph{Exact verification and remaining issue.}
We retain the rational exterior-algebra program
\begin{center}
\texttt{check\_partial\_fourier\_parity.py}.
\end{center}
It checks (PF.5), (PF.7), the full diagonal Fourier identity,
4,913 integral parity cases, and 41 multiplication graphs using
rational exterior algebra. Its saved output records the actual
runtime. The proof of the universal identities is the exterior
calculation and the displayed formulas; finite testing is only
a check of their signs and implementation.
The remaining mathematical issue is the existence of an
integral algebraic correspondence with $t-s$ odd. Since the
converse above shows this issue is equivalent to the original
splitting assertion, the criterion is recorded as a
reformulation and a concrete method exclusion, not as movement
of the unresolved main gate.

For comparison, the new \emph{\'etale} integral projectors of
\citet[Theorem 1.7, Corollary 1.8, and \S B.6]{cyKahnDivided2026PF}
do not provide the required ordinary Chow correspondence.
The distinction between these two cycle groups is explicit
in that source, including Question 1.11. We have not replaced
the ordinary Chow group in condition (2) by its \'etale version.
